%% General-TBox ALC: a finite-type algorithm and combined complexity.

\subsection{General TBoxes and elimination of unrealizable types}
\label{sec:p09-alc-general-tboxes}

For a finite TBox \(\mathcal T\), a concept \(C\) is satisfiable
relative to \(\mathcal T\) when some \(I\models\mathcal T\) has
\(C^I\ne\varnothing\).
Subsumption \(C\sqsubseteq_{\mathcal T}D\) means
\(C^I\subseteq D^I\) in every such interpretation.
Knowledge-base satisfiability uses the full finite pair
\((\mathcal T,\mathcal A)\) from Definition~\ref{def:p09-alc}.
For instance checking, \(\mathcal K\models C(a)\) means that every
model of \(\mathcal K\) interprets \(a\) in \(C^I\).
The complexity below is combined complexity: the TBox, ABox, and
query are all part of the input.

\begin{hbdefinition}{def:p09-alc-complete-type}{Complete concept types}
Use the finite NNF closure \(\mathcal C\) of
Definition~\ref{def:p09-alc-tableau}, with \(m=|\mathcal C|\).
For an NNF concept \(E\), write \(\bar E\) for its NNF complement.
The complement swaps \(\top,\bot\), swaps \(A,\neg A\), swaps
\(\sqcap,\sqcup\) while complementing their operands, and satisfies
\[
 \overline{\exists r.E}=\forall r.\bar E,\qquad
 \overline{\forall r.E}=\exists r.\bar E.
\]
Choose the normalization by these recursive clauses, so
\(\overline{\bar E}=E\).
Induction on concepts gives
\((\bar E)^I=\Delta^I\setminus E^I\).
The closure contains each complement and each immediate sub-concept.

A type is a subset \(t\subseteq\mathcal C\) satisfying:
\[
\begin{array}{ll}
 \text{constants:}&\top\in t,\quad\bot\notin t;\\
 \text{complements:}&\text{exactly one of }E,\bar E\text{ belongs to }t
                         \text{ for every }E\in\mathcal C;\\
 \text{conjunctions:}&E\sqcap F\in t
                         \Longleftrightarrow E\in t\text{ and }F\in t;\\
 \text{disjunctions:}&E\sqcup F\in t
                         \Longleftrightarrow E\in t\text{ or }F\in t;\\
 \text{global inclusions:}&\mathsf{Glob}(\mathcal T)\subseteq t.
\end{array}
\]
The Boolean conditions apply whenever the displayed compound concept
belongs to \(\mathcal C\).
For types \(t,u\) and an input role \(r\), define compatibility by
\[
 t\leadsto_r u
 \quad\Longleftrightarrow\quad
 \forall(\forall r.E)\in t\quad E\in u.
\]
These finite sets describe possible truth values at one element.
They become realized types only after their existential requirements
can be met.
\index{description logic!type elimination}\index{concept type}
\end{hbdefinition}

\begin{hbmethod}{def:p09-alc-type-elimination}{Eliminating unrealizable types}
Let \(S_0\) be the set of all types from the preceding definition.
There are at most \(2^m\).
Define a decreasing sequence by
\[
 S_{i+1}=
 \left\{t\in S_i:
 \begin{array}{l}
  \text{for every }\exists r.E\in t,\text{ some }u\in S_i\\
  \text{satisfies }t\leadsto_r u\text{ and }E\in u
 \end{array}\right\}.
\]
Stop when \(S_{i+1}=S_i\), and call the resulting set \(S\).
Each strict change deletes a type, so this takes at most \(2^m\)
strict stages.

Accept the tested-concept/TBox/ABox input precisely when all the
following conditions hold.
First, \(S\ne\varnothing\).
Second, there is an assignment \(a\mapsto t_a\in S\) for the
individual names occurring in the ABox such that
\[
\begin{aligned}
 E(a)\in\mathcal A&\ \Longrightarrow\
                         \operatorname{nnf}(E)\in t_a,\\
 r(a,b)\in\mathcal A&\ \Longrightarrow\
                         t_a\leadsto_r t_b.
\end{aligned}
\]
Finally, if a concept \(C_0\) is being tested, some \(t\in S\)
contains \(\operatorname{nnf}(C_0)\).
The assignment condition is vacuous for an empty ABox, while
\(S\ne\varnothing\) enforces the nonempty-domain convention even
for a TBox-only input.
Different individual names may be assigned the same type.
\end{hbmethod}

\begin{hbtheorem}{thm:p09-alc-type-elimination}{proved=annex:p9.alc.2}{Correctness and bounds of \(\mathcal{ALC}\) type elimination}
The method accepts exactly the satisfiable inputs of
Theorem~\ref{thm:p09-alc-tableau}.
Every accepted input has a model with at most \(2^m\) elements.
If the total input length is \(N\), then \(m=O(N)\), and an explicit
deterministic implementation runs in time \(2^{O(N^2)}\).
With no ABox individual names, the displayed enumeration and elimination
run in time \(2^{O(N)}\).
\end{hbtheorem}

The model has one element for each surviving type.
It interprets a concept name by the types containing it and a role
by all compatible pairs of types.
Completeness of the type descriptions is essential to the truth lemma.
If \(\forall r.E\) is absent, its complement \(\exists r.\bar E\)
is present and requires a surviving counterexample successor.
If \(\exists r.E\) is absent, \(\forall r.\bar E\) excludes every
compatible \(E\)-successor.

\paragraph{Two effects of general inclusions.}
For
\[
 \mathcal T=\{A\sqsubseteq\exists r.B,\quad
              B\sqsubseteq\exists r.C,\quad C\sqsubseteq\bot\},
\]
no initial type contains \(C\).
Every type containing \(B\) is deleted in the first elimination stage,
because its required \(C\)-successor is unavailable.
All remaining \(A\)-types are then deleted because their \(B\)-successors
have disappeared.
The TBox itself is satisfiable, for example on a singleton with
\(A,B,C\) all empty, but the tested concept \(A\) is unsatisfiable.

For \(\mathcal T=\{A\sqsubseteq\forall r.B\}\) and
\[
 \mathcal A=\{A(a),\ r(a,b),\ \neg B(b)\},
\]
every permitted type of \(a\) contains \(\forall r.B\), and every
permitted type of \(b\) contains \(\neg B\).
There is no compatible pair for the asserted role edge.
The name-assignment test therefore rejects the knowledge base.

\paragraph{Cyclic inclusions and finite models.}
The TBox \(\{A\sqsubseteq\exists r.A\}\) has a model with a
nonempty \(A\): take one element in \(A\) and an \(r\)-loop.
The type algorithm permits precisely this reuse of a surviving
type as its own successor.
The blocking tableau and the type algorithm encode the reuse in
different finite search objects; their proofs establish the same
satisfiability condition.

\begin{hbtheorem}{thm:p09-alc-general-exptime}{proved=annex:p9.alc.3}{Combined complexity of general-TBox \(\mathcal{ALC}\)}
The following problems are \(\mathsf{EXPTIME}\)-complete under
polynomial-time many-one reductions:
concept satisfiability relative to a general finite TBox;
subsumption relative to such a TBox;
satisfiability of a finite \(\mathcal{ALC}\) knowledge base; and
concept-instance checking for a finite knowledge base.
The satisfiability classification also holds when an optional
tested concept accompanies the knowledge base.
The input has exactly the constructors, GCIs, and assertions of
Definition~\ref{def:p09-alc}.
\end{hbtheorem}

The upper bound follows from the type algorithm.
The lower bound in Annex~\ref{annex:p9.alc.3} encodes an alternating
polynomial-space computation by concept names and two role relations.
A binary counter decreases on every generated transition, ensuring that
an infinite successor chain cannot replace a finite accepting computation.
The other reasoning tasks then follow by explicit polynomial reductions.

\paragraph{Source presentations and algorithmic scope.}
The classification for concept satisfiability under general TBoxes is
Theorem~6 on draft p.~26 of Baader, Horrocks, and Sattler's chapter
\emph{Description Logics}.  The cited locators refer to its 47-page author
draft; the published chapter belongs to the 2008
\emph{Handbook of Knowledge Representation}
\cite[draft \S3.6.1, Proposition~5 and Theorem~6]{BaaderHorrocksSattler2008}.
That presentation obtains the upper bound with looping tree automata.
The finite-type algorithm and alternating-machine construction above
supply the proofs under the handbook's stated input convention.

Donini and Massacci give a separate tableau algorithm for a concept and
a general TBox.  Their Theorem~3.13 proves single deterministic
exponential time for strategies satisfying their prescribed search
techniques; Theorem~6.17 supplies the depth-first-search call bound
\cite[pp.~98, 123]{DoniniMassacci2000}.
Their input called a knowledge base consists of the TBox.
The finite ABox assignment in the present method is an additional
construction proved in Annex~\ref{annex:p9.alc.2}.
The termination theorem for the earlier blocking tableau makes no
runtime claim about arbitrary expansion strategies.
